\documentclass[11pt,a4paper]{article}
\usepackage[margin=25mm]{geometry}
\usepackage{amsmath,amssymb,amsthm,booktabs,longtable,hyperref}
\setlength{\emergencystretch}{2em}
\newtheorem{theorem}{Theorem}\newtheorem{proposition}[theorem]{Proposition}
\newcommand{\gen}[1]{\langle#1\rangle}\newcommand{\Aut}{\operatorname{Aut}}
\title{Highly symmetric hypertopes: distinguished groups and finite constructions}
\author{Human Compression project}\date{}
\begin{document}\maketitle
\begin{abstract}
We assemble the general and finite results of Fernandes--Leemans--Weiss in
their fixed accepted manuscript. Residue stabilizers supply the C-group and
C$^+$-group assertions; an explicit chamber reflection supplies the converse
under its geometric hypotheses. Complete certificates identify all twenty-three
finite table entries and the three examples. The synthesis retains the
classical string and toroidal classification inputs, corrects the finite
captions and one group name, and separates the source's open questions.
The two complete component manuscripts and their certificates are part of
this proof, rather than replaced by the present summary.
\end{abstract}

\section{The principal statements and their hypotheses}
A hypertope is a thin, residually connected incidence geometry of finite
rank. Every flag extends to a chamber. Automorphisms in this paper preserve
types; correlations that permute types are not included in $\Aut_I$.
A regular hypertope has a chamber-transitive automorphism group. A chiral
hypertope has exactly two chamber orbits, with adjacent chambers in
different orbits. Propagation through residues makes both actions free.

Write $I=\{0,\ldots,r-1\}$. For involutions $\rho_i$ put
$G_S=\gen{\rho_i:i\in S}$. A C-group satisfies
$G_S\cap G_T=G_{S\cap T}$ for every $S,T\subseteq I$.
For generators $\alpha_1,\ldots,\alpha_{r-1}$ put $\alpha_0=1$ and
\[
 H_S=\gen{\alpha_i^{-1}\alpha_j:i,j\in S}.
\]
The C$^+$ condition is $H_S\cap H_T=H_{S\cap T}$ for
$|S|,|T|\ge2$, with trivial groups on singleton and empty supports.

\begin{theorem}\label{principal}
The principal results in the fixed source hold with the following precise
conventions.
\begin{enumerate}
\item A regular hypertope has distinguished nonidentity involutions
$\rho_i$ forming a C-group. Conversely, the standard parabolic coset system
of a C-group is a regular hypertope if the designated group is flag-transitive.
Its full type-preserving automorphism group is then that group.
\item A chiral hypertope of rank at least three has independent distinguished
generators $\alpha_i$ forming a C$^+$-group. Its group is transitive on flags
of every proper type, so every proper truncation is flag-transitive;
thinness of a truncation is not asserted.
\item Let a C$^+$-group $G$ define the standard coset system with
$P_i=H_{I\setminus\{i\}}$. Assume that this system is a hypertope and that
$G$ has exactly two chamber orbits. It is chiral, with full group $G$, if
and only if no automorphism of $G$ simultaneously inverts all the
$\alpha_i$. Otherwise its full group is $G\rtimes\gen F$, where $F$ is
a geometric involution interchanging the chamber orbits. If $\theta$ is
the inverting group automorphism, $FgF^{-1}=\theta(g)$.
\item At each of the twelve parameters of source Table~1 and eleven
parameters of source Table~2, the specified presentations give the finite
regular or chiral objects and full group orders in Table~\ref{finite} below,
with the corrected group identification stated there.
\end{enumerate}
\end{theorem}
\begin{proof}
We acquire the complete general proof~\cite{general} and finite
proof~\cite{finite}, together with their distinct reviews. The essential
interfaces are as follows. Fix a chamber $C=\{c_i:i\in I\}$ and let
$F_{I\setminus S}$ be its complementary flag. With left actions,
$C^{ij}=(C^i)^j$, the general proof establishes the identities
\[
 G_S=\operatorname{Stab}_G(F_{I\setminus S}),\qquad
 H_S=\operatorname{Stab}_G(F_{I\setminus S})
\]
in the regular and chiral cases respectively. In the second identity,
$\alpha_iC=C^{0i}$ and
$\alpha_i^{-1}\alpha_jC=C^{ij}$. Actual paths within the residue prove
the first identity, and even paths prove the second; neither proof assumes
an unidentified residue group. Intersections of the stabilizers give both
intersection conditions and the independence assertion. The standard
C-group converse uses flag transitivity to identify every flag and its
residue, and then the original intersection property proves thinness and
residue connectivity.

For the third assertion, the same proof identifies \emph{every} chamber
as $(g,+)=gC$ or $(g,-)=gC^0$, with
\[
                (g,+)\sim_i(g\alpha_i^{-1},-).
\]
The acquired two-orbit theorem gives the reflection
$(g,\pm)\mapsto(\theta(g),\mp)$. Reconstruction from colored chamber
adjacencies makes it an incidence automorphism. This is a geometric
involution even if $\theta$ is the identity. An abstract semidirect
product, without this reconstruction, would not pay the converse.

The finite component checks the hypotheses just used for all twenty-three
presentations: their exact finite groups, all marked intersections, actual
flag completions and residue graphs, and all-relator inversion tests.
Its regular completions and named-group maps are explicit isomorphisms.
For Table~1 its forty-eight subgroup-product inclusions additionally prove
that incidence ordered by type is transitive, giving the full polytope
axioms. Its computations use right cosets; the bijection
$P_i g\mapsto g^{-1}P_i$ preserves types and incidence and converts them
to the left-coset convention of the general proof. Thus these are the same
geometries. No new finite search is required for this join.
\end{proof}

\section{The finite rows and source corrections}
The source's printed $g$ is the order of the full type-preserving group.
It is the number of chambers in a regular row, and \emph{half} that number
in a chiral row. Thus both captions require the same correction.
\begin{center}
\begin{tabular}{ccrrc}\toprule
$p$&$(a,b)$&Table 1 group order&Table 2 group order&status\\\midrule
3&(2,0)&240&120&regular\\
3&(3,0)&1296&648&regular\\
3&(4,0)&15360&7680&regular\\
3&(1,2)&336&168&chiral\\
3&(1,3)&2184&1092&chiral\\
3&(1,4)&8064&4032&chiral\\
3&(2,2)&2880&1440&regular\\
3&(2,3)&6840&3420&chiral\\
4&(1,1)&288&---&regular\\
4&(2,0)&768&384&regular\\
4&(1,2)&2016&1008&chiral\\
5&(2,0)&28800&14400&regular\\\bottomrule
\end{tabular}
\end{center}
\refstepcounter{table}\label{finite}
\begin{center}\small Table \thetable. Exact positive finite roster; no completeness claim for chiral rows.\end{center}
The Table~2 group at $(3,2,2)$ is
\[
 D=\{(s,t)\in S_5\times S_4:\operatorname{sgn}(s)=\operatorname{sgn}(t)\},
\]
not $A_5\times S_4$. The finite proof gives a bijective map to $D$ and
distinguishes the two groups by the centralizers of their characteristic
perfect residuals, of orders $12$ and $24$. It identifies the groups
abbreviated $[1\,1\,2]^3$ and $[1\,1\,2]^4$ by explicit $G(3,3,4)$ and
$G_{29}$ models. Every other printed group name is likewise identified by
an actual map, including the actions in the split extensions. Each of the
eleven Table~2 full groups embeds with index two in its corresponding
Table~1 group; the complete marked halving maps and complements are retained.

The three examples are also acquired in full. Example~4.4 has group
$(C_3\times C_3)\rtimes C_2$ and incidence graph $K_{3,3,3}$, with
$27$ chambers and full group $S_3^3$. Its designated order-$18$ group is
not flag-transitive and its panels have three completions. Example~4.5
has the printed C-marking of $A_6$, $720$ chambers, $360$ panels with four
completions and $36$ disconnected residues. Example~6.3 has order $72$,
satisfies the involutory intersection property, and has rotation group
equal to the whole group; two rotation subgroups intersect in order $8$
instead of the required $4$. These are separate reasons why the relevant
geometric or index-two hypotheses cannot be discarded.

One additional source assertion about Example~4.5 needs qualification.
\begin{proposition}\label{a6}
The rotation marking of Example~4.5 is not a C$^+$-marking. Consequently
the sentence on source page~17 does not supply its claimed C$^+$ example.
\end{proposition}
\begin{proof}
Retain the source involutions $\rho_0,\rho_1,\rho_2,\rho_3$ and set
$\alpha_i=\rho_0\rho_i$. The acquired finite certificate gives $G=A_6$, the two parabolics
$G_{\{0,2,3\}}$, $G_{\{0,1,3\}}$ of order $60$, and their intersection
$G_{\{0,3\}}$ of order $10$. Every index-six subgroup of $A_6$ is
isomorphic to $A_5$: the action on its six cosets is faithful by simplicity,
has image $A_6$ in $S_6$, and identifies it with a point stabilizer.
For a group generated by involutions, the subgroup generated by their
pairwise products is normal and has index at most two. Thus simplicity
of $A_5$ gives
\[
 H_{\{0,2,3\}}=G_{\{0,2,3\}},\qquad
 H_{\{0,1,3\}}=G_{\{0,1,3\}}.
\]
Their intersection has order $10$, whereas
$H_{\{0,3\}}=\gen{\rho_0\rho_3}$ has order $5$.
Equivalently $\rho_0$ lies in the former intersection and not the latter.
This violates $IC^+$. The order and group-identification data are already
part of Example~4.5's accepted certificate; there is no additional census.
\end{proof}
This corrects the later illustration, not the local group and incidence
statements of Example~4.5. In particular it neither proves nor refutes a
general converse for arbitrary C$^+$-groups. Theorem~\ref{principal}(3)
keeps its hypertope and exactly-two-orbit hypotheses throughout.

\section{Complete source routing and retained inputs}
The fixed source is the twenty-two-page accepted author manuscript
of~\cite{source}. The complete route register accompanies this paper.
The nineteen named general clauses are routed as follows:
\begin{center}\small
\begin{tabular}{p{.25\textwidth}p{.66\textwidth}}\toprule
Source clauses&Complete acquired route\\\midrule
2.1--2.2&Propagation within every residue and flag extension.\\
3.1--3.5&Cosets and common-core quotient; actual chamber and panel stabilizers;
Dehon and Buekenhout--Hermand criteria.\\
4.1--4.3, 4.6&Regular residue identity, Tits' three-coset criterion,
rank-three panel completions and the flag-transitive converse.\\
5.1--5.2&Retained marked string correspondence and the proved rank-two
residue interpretation of a string diagram.\\
6.1--6.2&Independence for rotations and the even-word kernel at index two.\\
7.1--7.2&Even residue paths, full support intersections and independence.\\
8.1--8.2&Proper-flag transitivity and exact two-orbit coordinates.\\\bottomrule
\end{tabular}
\end{center}
Construction~8.1 defines an incidence system; flag extension and hypertope
properties are not automatic. A coset action is faithful only after factoring
its common core. In a nonfaithful chamber-transitive hypertope action the
Borel subgroup equals that core, rather than necessarily being trivial.
The proof of source Proposition~2.1 uses connectedness of each residue in
its induction, not connectedness of the whole incidence graph alone.
These qualifications are explicit in the general component.

Four previously established inputs remain attributed premises, rather than
new classification proofs: the marked string C-group correspondence
(Aschbacher, Schulte and McMullen--Schulte); the rank-three toroidal chirality
criterion of Breda D'Azevedo--Nedela; the toroidal hypermap families of
Corn--Singerman; and McMullen--Schulte's finite regular locally toroidal
polytope classification at $p=3,4,5$.
The toroidal inputs retain their admissible, nondegenerate parameter domains;
the unrestricted wording on source pages~12--13 cannot include the unit
period $(1,0)$. Indeed its displayed rotation presentation gives $y=x^2$
and the cyclic group $C_3$, with coincident parabolics and no thin geometry.
The exceptional $(1,1)$ system is Example~4.4. We do not assert a larger
parameter domain by extrapolating the finite tables. The positive rows above
are verified independently of classification completeness.

Earlier attributed remarks about polytope dualities, the $11$-cell,
Schulte--Weiss rotations, chiral extensions, codimension-two faces and the
Coxeter honeycomb construction remain contextual antecedents. They are not
new classifications, additional finite table rows or hypotheses secretly
used in Theorem~\ref{principal}. The $p=6$ programme is explicitly deferred
by this source; this synthesis does not turn that announcement into a new
L052 obligation or a fresh intake.

\section{Questions, evidence and attribution}
Problems~9.1 and~9.2 and the unnumbered completeness conjecture for the
known chiral list remain questions. Checking the positive rows does not
prove the latter two classifications. The programme's earlier negative
answers to Problems~9.3 and~9.4 are acquired with their precise standard
parabolic hypotheses: in~9.3 thinness and chamber connectedness force the
designated chamber orbit to be all chambers; in~9.4 the rank-three standard
construction is residually connected. The latter uses a connected spanning
coset graph, not an unjustified equality of incidence relations. These old
answers receive no new question credit and retain their previous limited
formalization scope.

The mathematical results and classical inputs keep the attributions in
\cite{source}. The general component was authored by the root collaborator
and independently reviewed by the present assembler. The finite component
was authored by a separate collaborator and independently reviewed by root.
The present contribution is the complete source join and
Proposition~\ref{a6}; the finite author confirmed its convention without
supplying a new finite computation. The accepted chamber and L003 shared
proofs remain complete acquired dependencies. Named group constructors,
exact presentation/isomorphism routines and the named CHEVIE model are
explicit software and identification premises of the finite certificate.
No new Lean theorem, measurement, compression factor or priority claim is
made. Whole-source admission still requires a distinct review of this
assembly and the publication gate.

\begin{thebibliography}{9}\small
\bibitem{source} M. E. Fernandes, D. Leemans and A. Ivi\'c Weiss,
\emph{Highly symmetric hypertopes}, Aequationes Math. 90 (2016),
1045--1067; fixed accepted author manuscript, 22 pages.
\bibitem{general} \emph{Residue stabilizers and distinguished generators},
accepted complete general recovery, eight pages; exact author, independent
review and root-adoption editions acquired.
\bibitem{finite} \emph{Finite coset geometries in Highly symmetric hypertopes},
accepted complete finite recovery, five pages; exact author and distinct
root-review editions acquired.
\bibitem{MS} P. McMullen and E. Schulte, \emph{Abstract Regular Polytopes},
Cambridge University Press, 2002. Marked correspondence and the cited regular
locally toroidal classifications retained.
\bibitem{BN} A. Breda D'Azevedo and R. Nedela, \emph{Chiral hypermaps of small
genus}, Beitr. Algebra Geom. 44 (2003), 127--143. Retained source input.
\bibitem{CS} D. Corn and D. Singerman, \emph{Regular Hypermaps},
J. Combinatorics 9 (1988), 337--351, as cited in the fixed source.
\end{thebibliography}
\end{document}
